\section{问题三：最小螺距的确定}
\subsection{建模思路}
问题三要求确定龙头前把手能够盘入半径为4.5 m的调头空间边界时所需的最小螺距。螺距越小，相邻圈螺线越密，龙头附近的板凳越容易与外圈板凳接近。因此，我们先以题目给定的55 cm作为初始螺距，并逐步减小螺距进行分析。

当螺距为55 cm时，根据前两问的计算过程，舞龙队在到达调头空间边界之前仍具有较大的空间余量。因此，最小螺距必然不大于55 cm。接下来从该上界出发，按步长逐步缩小螺距；每给定一个螺距，就判断龙头前把手从初始位置盘入到调头空间边界的过程中是否出现接触风险。若没有出现接触风险，则继续减小螺距；若出现接触风险，则退回上一螺距，并缩小步长继续搜索，直到步长达到给定精度。

\subsection{建立模型}
沿用问题一中的等距螺线模型，设螺距为 \(d\)，则盘入螺线可表示为
\[
r(\theta)=\frac{d}{2\pi}\theta,\qquad
x(\theta)=r(\theta)\cos\theta,\qquad
y(\theta)=r(\theta)\sin\theta.
\]
当龙头前把手到达调头空间边界时，有
\[
r(\theta_0)=R,\qquad R=4.5\ \text{m},
\]
即
\[
\theta_0(d)=\frac{2\pi R}{d}.
\]

龙头前把手到达调头空间边界时的几何关系示意图如下。图中边界圆与盘入螺线的交点对应$\theta_0(d)$，也是本问判断“能否盘入调头空间”的终点状态。

\begin{figure}[H]
\centering
\placeholderimage{5.4cm}{
图位预留：龙头到达调头边界示意图。\\[0.4em]
图中标出半径$R=4.5$ m的调头圆、盘入螺线$r=d\theta/(2\pi)$、\\
边界交点以及龙头前把手到达边界时的参数$\theta_0(d)$。
}
\caption{龙头前把手到达调头空间边界的几何示意}
\label{fig:q3-boundary}
\end{figure}

对于给定的 \(\theta\) 与 \(d\)，由相邻把手距离递推得到前若干把手坐标
\[
P_k(\theta,d)=\bigl(x_k(\theta,d),y_k(\theta,d)\bigr),\qquad k=0,1,\ldots,35.
\]
其中，龙头板凳的把手间距取 \(\ell_H=2.86\) m，其余板凳的把手间距取 \(\ell_B=1.65\) m。递推时通过
\[
\|P_{k+1}(\theta,d)-P_k(\theta,d)\|_2=\ell_k
\]
确定下一把手对应的角参数。

在临界碰撞判断中，主要考察龙头及其后一节板凳与前方若干龙身板凳之间的相对位置。对第 \(j\) 块板凳，其中心轴线由相邻把手 \(P_j,P_{j+1}\) 确定，记为
\[
A_jx+B_jy+C_j=0,
\]
其中
\[
A_j=y_{j+1}-y_j,\qquad
B_j=x_j-x_{j+1},\qquad
C_j=-x_jy_{j+1}+x_{j+1}y_j.
\]

龙头和第1节龙身的外侧角点由把手连线方向、法向方向以及板凳宽度确定。设这些角点构成集合 \(\mathcal{V}(\theta,d)\)，角点 \(V=(x_V,y_V)\) 到第 \(j\) 块板凳轴线的距离为
\[
\rho(V,L_j)=
\frac{|A_jx_V+B_jy_V+C_j|}
{\sqrt{A_j^2+B_j^2}}.
\]
取 \(j=3,\ldots,30\) 作为待检板凳范围，定义
\[
D(\theta,d)=
\min_{\substack{V\in\mathcal{V}(\theta,d)\\3\le j\le 30}}
\rho(V,L_j).
\]

由于板凳宽度为0.30 m，半宽为0.15 m。当龙头附近外侧角点到其他板凳轴线的最小距离降至0.15 m时，可认为实体轮廓到达接触临界。因此，判定函数为
\[
F(d)=\min_{\theta\in[\theta_0(d),\,\theta_0(d)+6]}D(\theta,d)-0.15.
\]
若 \(F(d)>0\)，说明在该螺距下未达到接触临界；若 \(F(d)\le 0\)，说明该螺距下已经出现接触风险。

\subsection{搜索最小螺距}
为了搜索最小螺距，我们采用逐步逼近的方法。

\textbf{Step1：确定初始螺距。} 取题设螺距 \(d=0.55\) m作为初始值，并选取0.01 m作为初始步长。

\textbf{Step2：判断是否发生碰撞。} 对当前螺距，计算龙头到达调头空间边界时的 \(\theta_0(d)\)，并在 \([\theta_0(d),\theta_0(d)+6]\) 内检查龙头及第1节龙身与第3至第30节龙身之间的距离关系。

\textbf{Step3：循环迭代。} 若未达到接触临界，则继续减小螺距；若达到接触临界，则返回上一螺距，并将步长缩小为原来的十分之一，继续判断。

\textbf{Step4：终止。} 当步长缩小到 \(10^{-6}\) m时停止循环，得到较为精确的最小螺距。

根据上述过程，得到不同搜索步长下的最小螺距近似值，如表\ref{tab:q3-result}所示。

\begin{table}[H]
\centering
\caption{逐步逼近最小螺距}
\label{tab:q3-result}
\begin{tabularx}{0.82\textwidth}{CCC}
\toprule
步长 & 螺距 & 单位\\
\midrule
0.01 & 0.46 & m\\
0.001 & 0.451 & m\\
0.0001 & 0.4504 & m\\
0.00001 & 0.45034 & m\\
0.000001 & 0.450332 & m\\
\bottomrule
\end{tabularx}
\end{table}

保留两位小数后，盘入调头空间所需的最小螺距为
\[
d_{\min}=0.45\ \text{m}.
\]
在粗搜索阶段，螺距 \(d=0.45\) m 已触发接触风险，这说明最小可行螺距位于0.45 m的右侧；继续细化后得到上述结果。
